\chapter{Improper integrals and tails}
\SourceRef{Zorich I, Section 6.5, PDF pp. 413--421; Apostol, Section 10.13; Zorich II, Section 17.2 for the quantitative tail form}
All improper integrals below require proper Riemann integrability on each
regular compact truncation. Infinite endpoints and finite singular endpoints
are treated by limits, not by assigning an integrand value at a singularity.

\KEntry{K03.10.one_endpoint}{One singular endpoint}
\begin{definition}
Let $b\in\R\cup\{+\infty\}$ and $a<b$. If $f$ is properly integrable
on every $[a,t]$, $a<t<b$, define
\[\int_a^bf=\lim_{t\to b-}\int_a^tf
\]
only when this limit is finite in $\R$ or $\C$. For $b=+\infty$ the
notation means $t\to+\infty$. Define a singular left endpoint by the
corresponding right-hand limit, with orientation as in module 08.
Infinite divergence is a separate conclusion, not convergence to a scalar.
\end{definition}

\KEntry{K03.10.several_endpoints}{Independent singular endpoints}
\begin{definition}
For $f$ locally integrable on $(a,b)$, choose $c\in(a,b)$ and define
$\int_a^bf=\int_a^cf+\int_c^bf$ only if both terms converge
independently. If there are finitely many interior singularities, split
at each and require both adjacent one-sided integrals separately.
\end{definition}
\begin{proposition}
Existence and value are independent of all regular splitting points.
\end{proposition}
\begin{proof}
For $c<d$ in a regular region and a left truncation $u<c$,
$\int_u^df=\int_u^cf+\int_c^df$. The last term is a fixed proper
integral, so the limit for one splitting point exists if and only if
the other does, with their values differing by that proper integral.
The same calculation for right truncations cancels this difference in
the total. Repeat for every regular component. This argument never
links two singular truncation parameters, so it does not authorize
cancellation between divergent endpoint terms.
\end{proof}

\KEntry{K03.10.cauchy}{The Cauchy tail criterion}
\begin{theorem}
$\int_a^bf$ converges if and only if for every $\eps>0$ there is a
regular $r<b$ such that
\[\left|\int_u^vf\right|<\eps\qquad(r<u<v<b).
\]
At infinity replace $r<b$ by a sufficiently large $r$.
\end{theorem}
\begin{proof}
Write $H(t)=\int_a^tf$. Then $\int_u^vf=H(v)-H(u)$ by proper
additivity. A finite limit of $H$ makes these differences small.
Conversely choose any increasing cofinal sequence $t_n\to b$.
The tail condition makes $H(t_n)$ Cauchy, hence convergent to a scalar
$L$. Given $\eps$, choose a threshold for differences below $\eps/2$
and then a $t_n$ beyond it with $|H(t_n)-L|<\eps/2$. Every sufficiently
late $t$ has $|H(t)-H(t_n)|<\eps/2$ (reverse the order if necessary),
so $H(t)\to L$. Left endpoint versions follow by reversing orientation.
\end{proof}

\KEntry{K03.10.nonnegative}{Nonnegative improper integrals}
\begin{theorem}
For real locally integrable $f\ge0$ on $[a,b)$, the partial integrals
$H(t)=\int_a^tf$ increase. They converge finitely if and only if they
are bounded above; otherwise they tend to $+\infty$.
\end{theorem}
\begin{proof}
Order and additivity give $H(v)-H(u)=\int_u^vf\ge0$.
If the set of partial values has finite supremum $L$, choose a partial
value above $L-\eps$; all later values lie in $[L-\eps,L]$, proving
convergence. If unbounded, for each $A$ choose one value above $A$;
all later values stay above $A$, proving divergence to $+\infty$.
\end{proof}

\KEntry{K03.10.comparison}{Comparison and limit comparison}
\begin{theorem}
For eventually $0\le f\le g$, convergence of $\int g$ implies
convergence of $\int f$. If $f,g$ are eventually positive and
$f(x)/g(x)\to c\in(0,\infty)$ at the singular endpoint, their
improper integrals converge or diverge together.
\end{theorem}
\begin{proof}
Discard a regular initial segment, which contributes only a finite proper
value. Order bounds partial integrals of $f$ by those of $g$, so the
nonnegative boundedness criterion proves the first assertion. For the
second, sufficiently late points satisfy
$(c/2)g\le f\le(3c/2)g$. Apply comparison in both directions.
Eventual positivity is needed for this order argument; the ratio test
does not make the same conclusion for arbitrary sign-changing functions.
\end{proof}

\KEntry{K03.10.absolute}{Absolute and conditional convergence}
\begin{definition}
An improper integral is absolutely convergent if $\int|f|$ converges.
It is conditionally convergent if $\int f$ converges and $\int|f|$
does not. These definitions apply separately at every singular endpoint.
\end{definition}
\begin{theorem}
Absolute convergence implies convergence for real or complex scalars,
and $|\int f|\le\int|f|$.
\end{theorem}
\begin{proof}
On every regular compact interval $|f|$ is integrable by module 08 and
$|\int_u^vf|\le\int_u^v|f|$. Cauchy tails of the latter therefore
give Cauchy tails of the former. Scalar completeness proves convergence;
passing to the limit in the finite norm inequality proves the value bound.
The oscillatory example below proves that the converse fails.
\end{proof}

\KEntry{K03.10.p_test}{Power and logarithmic tests}
\begin{theorem}
For real $p$,
\[\int_1^\infty x^{-p}\dd x<\infty\iff p>1,\qquad
\int_0^1 x^{-p}\dd x<\infty\iff p<1.
\]
For $A>1$, $\int_A^\infty\dd x/(x(\log x)^q)$ converges exactly
when $q>1$.
\end{theorem}
\begin{proof}
For $p\ne1$ a truncation from $u$ to $v$ equals
$(v^{1-p}-u^{1-p})/(1-p)$, by the proved power derivative and
Newton--Leibniz. At infinity $v^{1-p}$ tends to zero for $p>1$
and to infinity for $p<1$; the resulting finite value for $p>1$ is
$1/(p-1)$. At zero $u^{1-p}$ tends to zero exactly for $p<1$,
giving $1/(1-p)$, and diverges for $p>1$.
For $p=1$, truncations are $\log v-\log u$, divergent at both
relevant endpoints. For the logarithmic test make the proper monotone
substitution $y=\log x$, $\dd y=\dd x/x$ on each $[A,R]$.
Since $\log R\to\infty$ this becomes the power test
$\int_{\log A}^\infty y^{-q}\dd y$. The starting point $A>1$
ensures positive logarithm and no extra singularity at $1$.
\end{proof}

\KEntry{K03.10.integral_test}{Series and integrals}
\begin{theorem}
If $f:[1,\infty)\to[0,\infty)$ is nonincreasing, then
$\sum_{n=1}^{\infty}f(n)$ converges if and only if $\int_1^\infty f$
does. If they converge, for every integer $N\ge1$,
\[\int_{N+1}^{\infty}f\le\sum_{n=N+1}^{\infty}f(n)
\le\int_N^{\infty}f.
\]
\end{theorem}
\begin{proof}
Monotone bounded functions are integrable on each finite interval.
For integer $n$, monotonicity gives
$f(n+1)\le\int_n^{n+1}f\le f(n)$.
Sum over consecutive unit intervals. Boundedness of the nonnegative
partial sums is equivalent to boundedness of partial integrals (an
arbitrary truncation lies between two neighboring integer truncations).
The respective monotone convergence criteria prove equivalence.
Summing the shifted inequalities from $N$ to a finite upper integer
and then taking its limit gives both stated remainder bounds.
\end{proof}

\KEntry{K03.10.dirichlet}{Dirichlet's oscillation test}
\begin{lemma}[Proper parts formula for an accumulation primitive]
If $f$ is Riemann integrable on $[r,s]$, $F(x)=\int_r^xf$, and
$g\in C^1([r,s])$, then
\[\int_r^sfg=F(s)g(s)-\int_r^sFg'.
\]
\end{lemma}
\begin{proof}
First suppose $g$ is real-valued; $f$ and $F$ may still be complex-valued. For a partition write $F_i=F(x_i)$, $g_i=g(x_i)$.
The difference between $\int fg$ and
$\sum_i g_{i-1}(F_i-F_{i-1})$ is at most
$\|f\|_\infty(s-r)\omega_g(|P|)\to0$.
Finite summation by parts writes the latter as
$F(s)g(s)-\sum_iF_i(g_i-g_{i-1})$, since $F(r)=0$.
By the mean value theorem $g_i-g_{i-1}=g'(\eta_i)\Delta x_i$.
Replace $F_i$ by $F(\eta_i)$ with error at most
$(s-r)\|g'\|_\infty\omega_F(|P|)\to0$.
The remaining sum tends to $\int Fg'$, proving the formula.
For complex-valued $g=u+iv$, apply the proved real-$g$ identity separately to $u$ and $v$ and combine the two identities by complex linearity of the proper integral. This supplies the complex case without a false exact complex mean-value tag.
This proof is necessary because $F'=f$ need not hold at discontinuities
of $f$; a merely integrable $f$ does not make $F$ a $C^1$ primitive.
\end{proof}
\begin{theorem}
Let $f$ be locally Riemann integrable on $[a,\infty)$ with bounded
accumulation $A(x)=\int_a^xf$, $|A(x)|\le C$. If real $g$ is $C^1$,
monotone and tends to zero, then $\int_a^\infty fg$ converges.
\end{theorem}
\begin{proof}
A monotone function tending to zero has constant sign, and $h=|g|$
is nonincreasing with $h\to0$. For $R<S$ the primitive
$F_R(x)=\int_R^xf=A(x)-A(R)$ satisfies $|F_R|\le2C$.
The lemma, with $h'\le0$, gives
\[\left|\int_R^Sfh\right|
\le2Ch(S)+2C\int_R^S(-h')=2Ch(R).
\]
The right side tends to zero uniformly for $S\ge R\to\infty$,
so the Cauchy criterion proves convergence of $fh$ and hence $fg$.
The proof also works at a finite singular endpoint with $R,S$ tending
to that endpoint in the specified direction.
\end{proof}

\KEntry{K03.10.abel}{Abel's test}
\begin{theorem}
If $\int_a^\infty f$ converges and real $g$ is bounded, $C^1$ and
monotone, then $\int_a^\infty fg$ converges.
\end{theorem}
\begin{proof}
Let $|g|\le B$. For given $\eps>0$, the Cauchy criterion gives
$|F_R(x)|=|\int_R^xf|<\eps$ for all $x\ge R$ and all sufficiently
large $R$. The proper accumulation parts formula gives
\[\left|\int_R^Sfg\right|
\le\eps\left(|g(S)|+\int_R^S|g'|\right)
=\eps(|g(S)|+|g(S)-g(R)|)\le3B\eps.
\]
This proves Cauchy tails. For $B=0$ the product is zero.
Bounded monotonicity also gives a finite limit $\ell$ of $g$; the
identity $fg=\ell f+f(g-\ell)$ explains the distinct limits, but the
displayed estimate proves convergence without an unproved interchange.
\end{proof}

\KEntry{K03.10.parts}{Improper integration by parts}
\begin{theorem}
On regular truncations let the proper parts formula hold for $u,v$.
If $u(t)v(t)$ has a finite singular-endpoint limit and either of
$\int uv'$ and $\int u'v$ converges, then the other converges and
\[\int_a^buv'=[uv]_a^b-\int_a^bu'v.
\]
For several singular endpoints require these limits separately at each.
\end{theorem}
\begin{proof}
For regular $t<b$, the proper formula is
$\int_a^tuv'=u(t)v(t)-u(a)v(a)-\int_a^tu'v$.
If the boundary product and one integral have finite limits, this
identity expresses the remaining partial integral as a difference of
convergent quantities and proves its finite limit and the formula.
Repeat separately at a singular left endpoint or on either side of an
interior singularity. If the product has no justified limit, subtracting
formal infinite boundary terms is not a valid application of this theorem.
\end{proof}

\KEntry{K03.10.substitution}{Improper substitutions}
\begin{theorem}
Suppose $\phi:[\alpha,\beta)\to[a,b)$ is a monotone $C^1$
parameterization with $\phi(\alpha)=a$, $\phi(t)\to b$ as
$t\to\beta-$, and is onto the indicated interval. If $f$ is locally
properly integrable, then
\[\int_a^bf=\int_\alpha^\beta f(\phi(t))\phi'(t)\dd t
\]
whenever either improper limit converges, and existence of one is
equivalent to existence of the other. Orientation handles decreasing
parameterizations and left endpoints.
\end{theorem}
\begin{proof}
On each regular $[\alpha,T]$ the monotone substitution theorem proves
integrability of the transformed product and gives
$\int_\alpha^T(f\circ\phi)\phi'=\int_a^{\phi(T)}f$.
Cofinality $\phi(T)\to b$ proves the transformed limit if the original
limit exists. Conversely, for each regular $r\in[a,b)$ choose
$T_r$ with $\phi(T_r)=r$, using surjectivity; as $r\to b$, such
$T_r$ leave every compact subinterval of $[\alpha,\beta)$, since its
image has maximum strictly below $b$. Thus $T_r\to\beta$ and the
same identity proves the original limit. Flat portions cause no problem.
A nonmonotone parameterization is not covered by this cofinal argument.
\end{proof}

\KEntry{K03.10.principal_value}{Symmetric principal values}
\begin{definition}
At an interior singularity $c$ define, when finite,
\[\operatorname{pv}\int_a^bf
=\lim_{\eps\downarrow0}\left(\int_a^{c-\eps}f+
\int_{c+\eps}^bf\right),
\]
assuming the other endpoints are regular or separately controlled.
On the whole line a symmetric value is
$\lim_{R\to\infty}\int_{-R}^Rf$, with any interior singularities
handled explicitly. The center and linked truncation rule are part of
the definition. Independent improper convergence implies the matching
principal value by addition of limits, but not conversely.
\end{definition}

\KEntry{K03.10.pv_counterexample}{A principal value without an improper integral}
\begin{example}
For $f(x)=1/x$ on $[-1,1]\setminus\{0\}$,
\[\int_{-1}^{-\eps}\frac{\dd x}{x}=\log\eps,\qquad
\int_\eps^1\frac{\dd x}{x}=-\log\eps.
\]
Their sum is zero for each $\eps$, so the symmetric principal value
is zero. The two terms separately tend to $-\infty$ and $+\infty$,
so the ordinary improper integral does not exist.
\end{example}

\KEntry{K03.10.oscillatory_example}{A conditionally convergent oscillatory integral}
\begin{example}
$\int_1^\infty\sin x\,\dd x/x$ converges, but not absolutely.
\end{example}
\begin{proof}
The primitive $\int_1^x\sin t\dd t=\cos1-\cos x$ is bounded by two,
and $1/x$ decreases differentiably to zero, so Dirichlet applies.
Since $\sin(\pi/2)=1$ and sine is continuous, choose
$0<\delta<\pi/2$ such that $\sin x\ge1/2$ on
$[\pi/2-\delta,\pi/2+\delta]$. For integer $k\ge1$, translate this
interval by $k\pi$ to obtain $I_k$; the sign-change identity for sine
gives $|\sin x|\ge1/2$ on every $I_k$. Thus
\[\int_{I_k}\frac{|\sin x|}{x}\dd x
\ge\frac{\delta}{(k+1)\pi}.
\]
The harmonic series diverges (group terms between consecutive powers of
two, each group having sum at least $1/2$). Partial integrals of the
absolute integrand therefore grow without bound. This proves conditional
convergence by an actual recurring-interval lower bound.
\end{proof}

\KEntry{K03.10.tail_bounds}{Quantitative oscillatory remainders}
\begin{theorem}
If $|A(x)|=|\int_a^xf|\le C$ and real $g\ge0$ is $C^1$,
nonincreasing and tends to zero, then for $S\ge R\ge a$,
\[\left|\int_R^Sfg\right|\le2Cg(R),\qquad
\left|\int_R^\infty fg\right|\le2Cg(R).
\]
\end{theorem}
\begin{proof}
The proof of Dirichlet already gives the first bound, with primitive
$A(x)-A(R)$ bounded by $2C$. Dirichlet proves the infinite integral
exists, so letting $S\to\infty$ proves the second bound by continuity
of modulus. If instead one assumes directly
$|\int_R^xf|\le C_R$ for all $x\ge R$, the same calculation improves
both bounds to $C_Rg(R)$. These estimates specify the tail, rather than
merely invoking convergence by name, and will be used uniformly in
parameters in module 14.
\end{proof}

\KEntry{K03.10.endpoint_counterexample}{Linked endpoint truncations are insufficient}
\begin{example}
On $(0,1)$ let $f(x)=1/(1-x)-1/x$. A primitive on this interval is
$H(x)=-\log(1-x)-\log x$. Since $H(1-\eps)=H(\eps)$,
$\int_\eps^{1-\eps}f=0$ for every $0<\eps<1/2$.
Nevertheless
$\int_\eps^{1/2}f=H(1/2)-H(\eps)\to-\infty$,
and
$\int_{1/2}^{1-\eps}f=H(1-\eps)-H(1/2)\to+\infty$.
Both independent endpoint conditions fail. A finite linked limit on two
moving endpoints therefore does not establish ordinary improper convergence.
\end{example}
